\chapter{Fixed Income and Markets}\label{ch:iv:fixed-income-and-markets}

``Tell me what happened in markets yesterday.'' One candidate lists five closing levels. Another says that two-year
yields rose twelve basis points after a strong payroll number, the curve flattened, the dollar rallied and equities
recovered in the afternoon, and then says what she would have expected to see in front-end futures. The interviewer
asked for the second answer; only one candidate knew that. Fixed-income questions test bond arithmetic done quickly
(price, yield, duration, convexity, forwards, carry) and a view of how markets connect; the second is examined through
the \omterm{def:iv:fixed-income-and-markets:recap}{market recap}. The instruments are One Quant Book~2's.

\section{Price, yield, duration and convexity at interview speed}

A bond's price change for a small yield change is $\Delta P/P \approx -D_{\mathrm{mod}}\,\Delta y + \tfrac12\mathcal
C(\Delta y)^2$, with modified duration $D_{\mathrm{mod}}$ and convexity $\mathcal C$; its DV01 is $D_{\mathrm{mod}} \times
P \times 10^{-4}$ per unit of notional (One Quant Book~2, chapter~3). Three numbers are worth remembering: a par bond's
modified duration is close to $(1 - (1+y)^{-n})/y$, about 8.1 for ten years at 4\%; its convexity is about 81; and the
convexity term is second order, a few basis points of price for a 100 basis point move.

\begin{example}[A 100 basis point sell-off]\label{ex:iv:fixed-income-and-markets:selloff}
A ten-year 4\% annual-coupon bond at par has $D_{\mathrm{mod}} \approx 8.11$ and $\mathcal C \approx 80.8$. If yields rise
to 5\%, duration alone predicts $-8.11\%$; with convexity, $-8.11\% + \tfrac12 \times 80.8 \times 0.0001 = -7.71\%$; the
exact repricing gives $-7.72\%$.
\end{example}

\section{Curves: forwards, carry and roll-down}

Zero rates imply forward rates: $(1 + z_2)^2 = (1 + z_1)(1 + f_{1,2})$ for annual compounding. A bond held for a period
earns its \emph{carry} (coupon income minus the cost of financing it, One Quant Book~1, chapter~7) and its
\emph{roll-down}: if the curve does not move, the bond's yield slides down an upward-sloping curve as its maturity
shortens, and its price rises.

\begin{omfigure}
\begin{tikzpicture}
\begin{axis}[width=10cm,height=5.2cm,xlabel={\small maturity (years)},ylabel={\small yield (\%)},
  xmin=0,xmax=30,ymin=2.8,ymax=4.3,
  tick label style={font=\footnotesize},grid=major,grid style={gray!20},table/col sep=comma]
\addplot[omDef,thick] table[x=t,y=y]{figdata/interviews/18-fixed-income-and-markets/curve.csv};
\addplot[only marks,mark=*,mark size=2pt,omThm] table[x=t,y=y]{figdata/interviews/18-fixed-income-and-markets/points.csv};
\draw[-{Stealth},thick,omThm] (axis cs:5.1,3.76) to[bend right=40] (axis cs:3.9,3.68);
\draw[gray,thin] (axis cs:6.2,3.47) -- (axis cs:4.2,3.58);
\node[font=\scriptsize,anchor=north west,align=left] at (axis cs:6,3.5) {a five-year bond, one year later,\\is a four-year bond: its yield\\falls about 9 bp};
\end{axis}
\end{tikzpicture}

\omcaption{The chapter's synthetic upward-sloping curve (a Nelson--Siegel shape rising from about 2.9\% to 4.1\%). A
five-year bond held for a year rolls to the four-year point; with the curve unchanged its yield falls from 3.69\% to
3.59\%, which adds about 0.33\% to its price. Data: \texttt{fig\_iv\_rolldown.py}.}\label{fig:iv:fixed-income-and-markets:roll}
\end{omfigure}

\begin{method}[Carry and roll-down of a bond position]\label{met:iv:fixed-income-and-markets:carry}
\begin{enumerate}
\item Carry: coupon income over the horizon minus the repo cost of the price paid.
\item Roll-down: reprice the bond at the curve's yield for its shorter maturity at the horizon, curve unchanged.
\item The sum is the return if nothing happens; compare it with the yield move that would wipe it out, which is the
return divided by the duration at the horizon.
\end{enumerate}
\end{method}

\section{Repo, swaps and the basis in one paragraph each}

\emph{Repo.} A bond bought with borrowed cash, the bond pledged as collateral, earns the coupon and pays the repo rate on
the price: a positive carry when the coupon yield exceeds the repo rate, negative when it does not (One Quant Book~2,
chapter~5). \emph{Swaps.} A par swap rate is the fixed rate that makes an interest-rate swap worth zero; the swap spread,
swap rate minus the government yield of the same maturity, can be negative at long maturities when holding government bonds
uses scarce balance sheet (One Quant Book~2, chapter~9). \emph{FX forwards.} Covered interest parity makes the forward
$F = S(1 + r_d T)/(1 + r_f T)$ with simple rates, $S$ in quote currency per unit of base; deviations are the cross-currency
basis (One Quant Book~2, chapter~16).

\section{The market recap}

\begin{definition}[Market recap]\label{def:iv:fixed-income-and-markets:recap}
A \emph{market recap}\index{market recap} is a candidate's short account of a trading session asked for in an interview:
the level and change of a few key prices, the cause, the links across asset classes, and what the candidate would watch
next, in two minutes.
\end{definition}

The recap is scored on structure and on causal links, not on recall. Book 2, chapter~31 reads a macro calendar and a rates
screen; the method below turns that reading into an answer.

\begin{method}[Giving a market recap]\label{met:iv:fixed-income-and-markets:recap}
\begin{enumerate}
\item Lead with the day's driver: an economic release against its consensus forecast, a central-bank statement, an
event.
\item Give the rates move by maturity and what it did to the curve (steepener or flattener), in basis points.
\item Link: the currency, equity indices, credit spreads, volatility, commodities, each with a direction and a reason.
\item Say what surprised you (a move that did not fit) and what you would watch tomorrow.
\item Use numbers, rounded; never read a list of closing levels.
\end{enumerate}
\end{method}

\begin{example}[A recap from a table]\label{ex:iv:fixed-income-and-markets:recap}
A synthetic session: a payroll number well above consensus; two-year yields $+12$ bp, ten-year $+5$ bp; the dollar $+0.8\%$;
an equity index $-0.9\%$ at the open and $-0.4\%$ at the close. Recap: ``A strong payroll print moved front-end rate
expectations up, so two-year yields rose 12 basis points and the two-to-ten curve flattened 7; higher expected rates lifted
the dollar. Equities fell at the open on rates and recovered half of it as the growth reading was digested. I would watch
whether the front end holds the move into the next inflation release.''
\end{example}

\section{Worked answers}

\begin{example}[The DV01 of a zero-coupon bond]\label{ex:iv:fixed-income-and-markets:zero}
\emph{``A ten-year zero-coupon bond yields 4\% with annual compounding. What are its price and its DV01 per 100 of face?''}
The price is $100/1.04^{10}$; $1.04^{10} \approx e^{0.392} \approx 1.480$, so about 67.56. A zero's Macaulay duration is
its maturity, 10, and its modified duration $10/1.04 \approx 9.62$; the DV01 is price times modified duration times one
basis point, $67.56 \times 9.62 \times 10^{-4} \approx 0.065$. \emph{Check:} repricing at 4.01\% gives a fall of 0.0649,
the same to the precision quoted; the small gap is convexity.
\end{example}

\begin{example}[A breakeven inflation rate]\label{ex:iv:fixed-income-and-markets:breakeven}
\emph{``A ten-year nominal government bond yields 4.2\% and the inflation-linked bond of the same maturity has a real
yield of 1.9\%. What inflation does the market price, and what else is in that number?''} To first order, the breakeven is
the difference, 2.3\%; exactly, with annual compounding, $1.042/1.019 - 1 \approx 2.26\%$. Then the second half, which the
interviewer is waiting for: the breakeven is not a pure forecast. It contains an inflation risk premium (positive when
investors pay to be protected), a liquidity premium on the linked bond (which is less traded, lowering the breakeven),
and the index-lag and seasonality conventions of the linked bond. A candidate who says ``the market expects 2.3\%'' and
stops has given half an answer.
\end{example}

\section{Question bank}

\begin{interviewq}[$\star$ \iqroles{trader, bank} \iqfirm{bank}]\label{iq:iv:fixed-income-and-markets:1}
A bond has a modified duration of 7. Yields rise 25 basis points. What happens to its price, approximately?
\end{interviewq}

\begin{interviewq}[$\star$ \iqroles{trader, risk} \iqfirm{bank}]\label{iq:iv:fixed-income-and-markets:2}
What is the DV01 of a position of 10 million nominal in a bond priced at 98 with a modified duration of 6.5?
\end{interviewq}

\begin{interviewq}[$\star$ \iqroles{trader, researcher} \iqfirm{any}]\label{iq:iv:fixed-income-and-markets:3}
One-year and two-year zero rates are 3\% and 4\% (annual compounding). What is the one-year rate one year forward?
\end{interviewq}

\begin{interviewq}[$\star$ \iqroles{trader} \iqfirm{multi-manager fund}]\label{iq:iv:fixed-income-and-markets:4}
``Tell me what happened in markets yesterday.'' What structure does your answer follow, and what does the interviewer
not want to hear?
\end{interviewq}

\begin{interviewq}[$\star\star$ \iqroles{trader, risk} \iqfirm{bank}]\label{iq:iv:fixed-income-and-markets:5}
You are long 10 million of a bond with a DV01 of 850 per million. You hedge with a bond whose DV01 is 460 per million.
How much of the second bond do you sell? What risk remains?
\end{interviewq}

\begin{interviewq}[$\star\star$ \iqroles{bank, trader} \iqfirm{bank}]\label{iq:iv:fixed-income-and-markets:6}
A ten-year 4\% annual-coupon bond trades at par. Estimate its price change for a 100 basis point rise in yields with
duration alone and with convexity, and compare with the exact change.
\end{interviewq}

\begin{interviewq}[$\star\star$ \iqroles{trader, researcher} \iqfirm{multi-manager fund}]\label{iq:iv:fixed-income-and-markets:7}
On the chapter's curve, you buy the five-year bond with a coupon equal to its yield, financed at 3\% repo, and hold it for a
year with the curve unchanged. What are its carry, its roll-down and its total return?
\end{interviewq}

\begin{interviewq}[$\star\star$ \iqroles{trader} \iqfirm{bank}]\label{iq:iv:fixed-income-and-markets:8}
You buy 100 million of a bond with a 4\% coupon, priced at par, and finance it in repo at 4.5\% for one month. What is your
carry for the month? What must happen for the trade to make money?
\end{interviewq}

\begin{interviewq}[$\star\star$ \iqroles{trader} \iqfirm{bank}]\label{iq:iv:fixed-income-and-markets:9}
EURUSD spot is 1.10; one-year dollar and euro rates are 4\% and 2\% (simple). What is the one-year forward? If the market
forward is 1.1180, what does it tell you?
\end{interviewq}

\begin{interviewq}[$\star\star\star$ \iqroles{trader, bank} \iqfirm{bank}]\label{iq:iv:fixed-income-and-markets:10}
The thirty-year swap rate is below the thirty-year government yield. How can a swap rate be below the yield of a
government bond, and what would it take to earn the spread?
\end{interviewq}

\begin{interviewq}[$\star\star\star$ \iqroles{trader} \iqfirm{multi-manager fund}]\label{iq:iv:fixed-income-and-markets:11}
Give a two-minute recap of this synthetic session: a payroll number far above consensus; two-year $+12$ bp, ten-year $+5$
bp; the dollar $+0.8\%$; an equity index $-0.9\%$ at the open, $-0.4\%$ at the close; credit spreads unchanged. What would
have surprised you?
\end{interviewq}

\begin{interviewq}[$\star\star\star$ \iqroles{risk, researcher} \iqfirm{asset manager}]\label{iq:iv:fixed-income-and-markets:12}
Build a barbell of two-year and thirty-year par bonds (4\% coupons, 4\% yields) with the same modified duration as the
ten-year par bond. Compare their convexities. Which does better on a parallel move of 100 basis points either way, and
what is the catch?
\end{interviewq}

\begin{interviewq}[$\star\star\star$ \iqroles{trader, researcher} \iqfirm{any}]\label{iq:iv:fixed-income-and-markets:13}
A central bank raises rates by 50 basis points when 25 were expected. What do you expect for the front end, the long end,
the currency and equities, and why could each go the other way?
\end{interviewq}

\begin{omsources}
\item One Quant Book~1, chapter~7 (carry); One Quant Book~2, chapters 3, 5, 9, 16 and~31 (bonds, repo, swaps, FX forwards,
reading a macro calendar).
\item C.~R.~Nelson and A.~F.~Siegel, ``Parsimonious modeling of yield curves'', \emph{Journal of Business} 60(4), 1987 (the curve's
functional form).
\end{omsources}
